% Einheit 2 von 4 – Ableitung an einer Stelle (bank/ableitung-und-aenderungsrate/e2.jsonl, zusammenbau v0.1)
%% TODO Zeitmarke Einheit 2: Marken-Zeile nicht lesbar
%% TODO Zweigzeile Teil 1: was hier gelernt wird (Satz in Schülersprache aus dem Einheitstitel) – Einheit 2
\einheitenkopf[e2]{Einheit 2 von 4 · Ableitung an einer Stelle}
\zweigzeile{Abitur GK}
\setcounter{aufgabe}{11}

%% TODO Titel: Ich-kann-Satz fehlt in der Bank (Platzhalter: Kettenname) – Nr. 12
%% TODO Anweisung: ein Satz über der ersten Teilaufgabe fehlt in der Bank – Nr. 12
\begin{aufgabe}{Ableitung an einer Stelle}
%% TODO \rechenplatz{n}: Schrittzahl des Grundfalls fehlt in der Bank (nur bei fester Schreibform, 2.3 b)
\begin{teile}
\teil $h(t)$ ist die Höhe eines Baums in cm nach $t$ Jahren. Was ist $h'(4) = 35$? \\ \kreuz{Funktionswert – Höhe oder Bestand} \\ \kreuz{Ableitungswert – Steigung oder Rate}
\teil Die Höhe eines Balls wird durch $h(t) = -5t^2 + 20t$ beschrieben ($t$ in s, $h(t)$ in m). Berechne $h'(1)$ und gib an, was die Zahl bedeutet. h'(1) = \leerfeld[m] je s
\teil Die Abbildung zeigt den Graphen der linearen Funktion $g$ durch $(1 | -1)$ und $(3 | 5)$. Begründe, dass $g'(3) = 3$ gilt \hfill (Abitur 2026 GK).

\begin{ksys}[xmin=-1,xmax=4,ymin=-3,ymax=7,ablesen]\gerade{3}{-4}{g}\punkt{1}{-1}{}\punkt{3}{5}{}\end{ksys}
\teil Gegeben ist $g(x) = 3\mathrm{e}^x - 3$; die Abbildung zeigt den Graphen. Berechne $g'(0)$ und veranschauliche in der Abbildung, wie man diesen Wert grafisch ermittelt \hfill (Abitur 2024 GK).

\begin{ksys}[xmin=-3,xmax=2,ymin=-4,ymax=8,ablesen]\funktion{3*exp(\x)-3}{g}\end{ksys}
\teil Die Pollenanzahl je Kubikmeter wird durch $p(t) = 2t^2 - 40t + 300$ beschrieben ($t$ in Stunden nach Messbeginn, $0 \le t \le 10$). Zu welchem Zeitpunkt beträgt die momentane Änderungsrate $-16$ Pollen je Kubikmeter und Stunde? t = \leerfeld
\teil $f$ beschreibt die momentane Änderungsrate der Staulänge auf einer Autobahn in km je Stunde; $x$ ist die Zeit in Stunden nach 6:00 Uhr. Es gilt $f(3) < 0$. Gib die Bedeutung im Sachzusammenhang an \hfill (Abitur 2023 LK).
\teil $f$ beschreibt für $0 \le x \le 20$ die momentane Änderungsrate des Inhalts eines Tanks in kg je Stunde ($x$ Zeit in Stunden seit Beobachtungsbeginn); der Punkt $(3 | 180)$ liegt auf dem Graphen von $f$. Gib die Bedeutung der Koordinaten im Sachzusammenhang an \hfill (Abitur 2021 GK).
\teil Die Schaufel eines Wasserrads wird durch die Graphen von $f(x) = 0{,}1x^3 - x^2 + 2x$ und $p(x) = -x^2 + 3x - 0{,}5$ begrenzt. Weise nach, dass $p'(0{,}5) = f'(0)$ gilt, und erläutere die anschauliche Bedeutung im Sachzusammenhang \hfill (Abitur 2021 GK).
\teil Die Abbildung zeigt den Graphen einer Funktion $f$. Beurteile die Aussage: „Für $0 \le x \le 2$ ist die Steigung des Graphen kleiner als $2{,}5$.“ \hfill (Abitur 2023 LK)

\begin{ksys}[xmin=-1,xmax=3,ymin=-1,ymax=6,ablesen]\funktion{-\x^3+3*\x^2+1}{f}\end{ksys}
\end{teile}
\end{aufgabe}

%% TODO Titel: Ich-kann-Satz fehlt in der Bank (Platzhalter: Kettenname) – Nr. 13
%% TODO Anweisung: ein Satz über der ersten Teilaufgabe fehlt in der Bank – Nr. 13
\begin{aufgabe}{Momentane Änderungsrate am Graphen über eine eingezeichnete Tangente bestimmen}
\begin{teile}
\teil Die Abbildung zeigt das Wasservolumen $f(t)$ in einem Becken (in m³) in Abhängigkeit von der Zeit $t$ (in h). Bestimme die momentane Änderungsrate des Volumens zwei Stunden nach Beobachtungsbeginn, indem du die Tangente einzeichnest \hfill (Abitur 2017 LK).

\begin{ksys}[xmin=0,xmax=8,ymin=0,ymax=120,xstep=1,ystep=20,xlabel=t in h,ylabel=V in m³,ablesen]\funktion{5*\x^2}{f}\end{ksys}
\end{teile}
\end{aufgabe}

%% TODO Titel: Ich-kann-Satz fehlt in der Bank (Platzhalter: Kettenname) – Nr. 14
%% TODO Anweisung: ein Satz über der ersten Teilaufgabe fehlt in der Bank – Nr. 14
\begin{aufgabe}{Graph der Füllhöhe in Abhängigkeit von der Zeit über die Gefäßform auswählen und begründen}
\begin{teile}
\teil Eine nach oben weiter werdende Schale wird mit konstanter Zuflussrate gefüllt. Welcher Graph beschreibt die Füllhöhe in Abhängigkeit von der Zeit? Begründe \hfill (Abitur 2017 LK). \\ \kreuz{Graph I: steigt immer langsamer (rechtsgekrümmt)} \\ \kreuz{Graph II: steigt gleichmäßig (Gerade)} \\ \kreuz{Graph III: steigt immer schneller (linksgekrümmt)}
\end{teile}
\end{aufgabe}

%% TODO Titel: Ich-kann-Satz fehlt in der Bank (Platzhalter: Kettenname) – Nr. 15
%% TODO Anweisung: ein Satz über der ersten Teilaufgabe fehlt in der Bank – Nr. 15
\begin{aufgabe}{Ableitung an einer Stelle – Fehler finden, Begründen, Anwendung, Darstellungswechsel}
\begin{teile}
\teil Für $g(x) = 3\mathrm{e}^x - 3$ soll $g'(0)$ berechnet werden. Paul rechnet $g(0) = 3 - 3 = 0$ und antwortet: $g'(0) = 0$. Finde den Fehler.
\teil Begründe, warum die Ableitung einer linearen Funktion $g(x) = mx + n$ an jeder Stelle den Wert $m$ hat.
\teil Die Höhe eines Heißluftballons wird durch $h(t) = -0{,}5t^2 + 12t$ beschrieben ($t$ in Minuten nach dem Start, $h(t)$ in m). Der Pilot darf höchstens mit $4$ m je Minute steigen. Ab welchem Zeitpunkt hält er die Vorgabe ein?
\teil Für $f(x) = 0{,}5x^2$ ist $f'(2) = 2$. Zeichne die Tangente im Punkt $(2 | 2)$ mit einem passenden Steigungsdreieck in die Abbildung ein.

\begin{ksys}[xmin=-1,xmax=4,ymin=-1,ymax=6,ablesen]\parabel{0.5}{0}{0}{f}\end{ksys}
\end{teile}
\end{aufgabe}
